\section{La primera ola de modelos instruct abiertos: la distribución de los datos importa más que el nombre}

La sección anterior definió IFT; esta sección sigue a la primera ola de modelos de 2023. Alpaca demostró que la instruction data sintética permite adaptar LLaMA con rapidez; Vicuna demostró que la distribución de prompts de chat reales es decisiva; OpenAssistant demostró que los datos humanos abiertos tienen valor a largo plazo; StableVicuna mostró que PPO no es algo que solo las empresas de código cerrado puedan ejecutar. La variable central pasó de manera continua de «si existe el algoritmo» a «si se tienen los datos correctos».

\subsection{Alpaca e instruction fine-tuning}

Esta subsección comienza con Alpaca. Su importancia no está en su calidad absoluta medida con los estándares actuales, sino en que mostró lo siguiente: un base model sólido, más una cantidad modesta de instruction-response data, basta para obtener un comportamiento chat/instruct evidente. Al leer la figura, primero distinga el conocimiento del base model del formato de respuesta aprendido mediante fine-tuning.

\lecturefigure{slide-029.jpg}{Alpaca: uno de los primeros instruction-tuned model abiertos de amplia difusión}{official deck, p.~29}

\lecturefigure{slide-030.jpg}{Los dos objetivos de IFT: adaptarse al estilo de entrada y admitir system prompt y multi-turn}{official deck, p.~30}

Si la pérdida se calcula solo sobre los assistant answer token, el objetivo de IFT puede escribirse como
\[
\mathcal{L}_{\mathrm{IFT}}=-\sum_{t\in\mathcal{A}}\log p_{\theta}(y_t\mid s,u,y_{<t}),
\]
donde $s$ es el system prompt, $u$ es la user instruction, $y$ es la assistant response y $\mathcal{A}$ son las posiciones de assistant token que requieren supervisión. El chat template codifica los roles y los límites como token especiales; una plantilla que no coincide puede hacer que el mismo modelo rinda notablemente peor.

\lecturefigure{slide-031.jpg}{IFT parte del base model y continúa el entrenamiento con instruction-response pairs}{official deck, p.~31}

\begin{importantbox}{Qué cambia principalmente IFT}
IFT suele reforzar tres tipos de capacidad: reconocer que «esto es una solicitud», organizar la respuesta según el rol y el formato, y proyectar el conocimiento ya existente hacia la salida que pide la instrucción. Puede mejorar la capacidad en la tarea, pero también puede provocar el olvido o el desplazamiento de la base capability; por ello hay que evaluar a la vez el chat behavior y las capacidades originales, y no tomar «conversa mejor» como «es más fuerte en todos los aspectos».
\end{importantbox}

\teachervoice{00:13:42--00:14:53, Nathan describe IFT como el proceso que hace que el modelo integre el chat template, el system prompt, los roles en varios turnos y el estilo de entrada. Sigue usando una autoregressive loss; la diferencia proviene sobre todo de la estructura condicional de los ejemplos de entrenamiento.}

\subsection{Self-instruct: ampliar la instruction data con modelos}

La subsección anterior explicó que el cambio de comportamiento de IFT proviene principalmente de la distribución instruction-response; esta subsección plantea la pregunta siguiente: cuando los equipos abiertos no tienen presupuesto para anotación humana a gran escala, ¿de dónde sale esa distribución? La reproducibilidad de Alpaca proviene de self-instruct: se parte de unos pocos seed prompt escritos por personas, se pide a un modelo sólido que genere más instrucciones y respuestas, y después se filtran los ejemplos duplicados, de baja calidad o con errores de formato. Al leer la figura, considere la calidad de las seed, la diversidad de la generación y las reglas de filtrado como un mismo sistema de datos; este convierte la costosa redacción humana en un pipeline de «semillas humanas + ampliación con modelos + filtrado», y también introduce el riesgo de heredar los sesgos del modelo maestro y de contraer la distribución.

\lecturefigure{slide-032.jpg}{Self-instruct: ampliar la instruction data sintética a partir de pocas seed de alta calidad}{official deck, p.~32}

Sea $D_0$ el conjunto de seed, $G$ el generador y $F$ el filtro; una ronda de ampliación puede escribirse como
\[
D_{k+1}=D_k\cup F\bigl(G(D_k)\bigr).
\]
Aquí $G(D_k)$ genera nuevos pares instruction/response y $F$ elimina duplicados y aplica las reglas de calidad. Si $G$ y $F$ solo favorecen plantillas similares, aumentar el volumen de datos no aumentará la cobertura real de tareas.

\begin{warningbox}{Los dos riesgos de ciclo cerrado de la synthetic data}
Primero, el student hereda los errores factuales, los límites de rechazo y el estilo lingüístico del teacher; segundo, la self-generation repetida comprime la diversidad y hace que el modelo sea cada vez más hábil en «preguntas al estilo de un modelo». Por ello conviene mezclar datos de usuarios reales y datos humanos, y medir la cobertura semántica de prompt/response, no solo contar el número de ejemplos.
\end{warningbox}

\teachervoice{00:14:53--00:16:10, el curso considera self-instruct un factor clave en la explosión temprana de modelos abiertos: el modelo que genera las respuestas sustituyó parte de la anotación humana y volvió accesible la instruction data. En la segunda mitad de la clase se advierte además que la distribución sintética puede ser repetitiva y frágil.}

\subsection{Vicuna y ShareGPT: los prompt de usuarios reales cambian la distribución}

Self-instruct resuelve «cómo ampliar el volumen rápidamente», pero aún no responde «si las preguntas de entrenamiento se parecen a las que harían los usuarios reales». Por eso Vicuna desplazó la atención del número de ejemplos a la prompt source: ShareGPT reunió conversaciones de ChatGPT compartidas por usuarios, lo que permitió que un modelo abierto viera por primera vez solicitudes de varios turnos, preguntas de seguimiento y estilos más cercanos a un producto real. Al revisar las dos páginas siguientes, compare a la vez la synthetic instruction y las conversaciones reales en cuanto a longitud de la tarea, estructura de turnos y riesgos de privacidad; la mejora de calidad proviene de que la distribución coincide, y no solo de tener más token.

\lecturefigure{slide-033.jpg}{Vicuna, poco después de Alpaca, incorporó una chat prompt distribution más realista}{official deck, p.~33}

\lecturefigure{slide-034.jpg}{ShareGPT: origen de los datos, valor y riesgos de privacidad}{official deck, p.~34}

\begin{knowledgebox}{Lectura de la figura: el valor de los datos y la gobernanza de los datos aparecen juntos}
ShareGPT aportó lo que los usuarios reales quieren preguntar, una señal que la synthetic instruction difícilmente puede replicar; sin embargo, una herramienta para compartir no equivale por sí misma al consentimiento del usuario para usar sus datos en entrenamiento, y el texto puede contener además información personal. Los posteriores LMSYS-Chat-1M y WildChat dieron más importancia a la limpieza y al consent, lo que muestra que la legalidad y la calidad de los datos abiertos son requisitos del mismo pipeline.
\end{knowledgebox}

\teachervoice{00:17:21--00:18:55, Nathan dice que ShareGPT es una de las pocas fuentes que permitió a los builder abiertos ver prompt de usuarios reales, y señala que los conjuntos de datos posteriores dedicaron mucho trabajo a eliminar información personal o a obtener el consentimiento antes de la recopilación.}

\subsection{Weight differences y las restricciones para publicar en abierto}

Las dos subsecciones anteriores trataron cómo se obtienen los datos; esta pasa a otra restricción del sistema que suele pasarse por alto: terminar el entrenamiento no significa que el modelo pueda entregarse de forma legal y estable. La licencia original de los pesos de LLaMA y su forma de distribución hicieron incómoda la publicación de fine-tune: los equipos solo podían publicar la weight difference respecto del base model, y los usuarios debían obtener legalmente los base weights para después combinarlos. Al leer esta línea de tiempo, distinga la capacidad del modelo, la receta de entrenamiento y el mecanismo de distribución; la weight diff no es un algoritmo de aprendizaje automático, pero determina directamente si el modelo puede instalarse, reproducirse e iterarse.

\lecturefigure{slide-035.jpg}{Los modelos instruct abiertos siguen multiplicándose: Vicuna, Koala, Dolly, entre otros}{official deck, p.~35}

\lecturefigure{slide-036.jpg}{Por qué publicar weight differences: distribuir fine-tune bajo restricciones de licencia}{official deck, p.~36}

Si los pesos base son $W_0$ y los pesos fine-tuned son $W_1$, publicar la diferencia equivale a
\[
\Delta W=W_1-W_0,\qquad W_1=W_0+\Delta W.
\]
Aquí $\Delta W$ no reduce la cantidad de información de un fine-tune completo; solo traslada al usuario la dependencia de la distribución. Es distinto del update de bajo rango de LoRA: la weight diff puede ser un tensor de rango completo y del mismo tamaño.

\lecturefigure{slide-037.jpg}{La línea de tiempo de modelos incorpora Koala y Dolly, y la forma de publicación se estandariza gradualmente}{official deck, p.~37}

\lecturefigure{slide-038.jpg}{Evolución de las puntuaciones MT-Bench de los primeros modelos instruct abiertos}{official deck, p.~38}

\begin{warningbox}{Las puntuaciones históricas solo explican el ecosistema de su momento}
El judge, los prompt y las versiones de modelo de MT-Bench pueden haberse actualizado después. Las puntuaciones de la figura sirven para ilustrar los cambios relativos entre modelos de 2023; no pueden compararse directamente con modelos de 2026 ni con configuraciones de judge distintas. Al reproducirlas, registre el commit de la evaluación, el judge model, la temperature y el prompt template.
\end{warningbox}

\teachervoice{00:19:00--00:21:05, el ponente explica que la weight diff fue una respuesta de ingeniería a las restricciones de publicación de LLaMA, y que solo los cambios posteriores de licencia facilitaron la publicación directa. La velocidad del ecosistema abierto depende de la ley y de las interfaces de distribución, no solo del código de entrenamiento.}

\subsection{OpenAssistant y StableVicuna: datos humanos y PPO temprano}

Esta subsección examina dos líneas que suelen simplificarse. OpenAssistant publicó datos de conversación generados y anotados por personas, que durante mucho tiempo fueron materia prima para modelos posteriores; StableVicuna, por su parte, completó relativamente pronto un RLHF abierto con PPO. Ambos muestran que la comunidad abierta no se limitó a IFT sintético, sino que también exploró las preferencias humanas y los pipeline de RL de alto costo.

\lecturefigure{slide-039.jpg}{OpenAssistant Conversations: uno de los primeros conjuntos de datos abiertos de instruction humanas}{official deck, p.~39}

\lecturefigure{slide-040.jpg}{StableVicuna: uno de los primeros modelos RLHF abiertos en usar PPO}{official deck, p.~40}

\begin{importantbox}{Por qué la human data es difícil de sustituir}
Los datos humanos aportan intenciones reales, lenguaje diverso, casos límite y preferencias relativas; estas señales determinan cómo el modelo ordena respuestas que son «razonables, pero no únicas». El costo proviene del reclutamiento, la capacitación, la consistencia, la privacidad y el control de calidad; por ello, lo más escaso para la comunidad abierta no suele ser el script de entrenamiento, sino datos humanos de alta calidad que puedan publicarse legalmente.
\end{importantbox}

\teachervoice{00:23:43--00:24:47, Nathan considera OpenAssistant uno de los proyectos abiertos tempranos más exitosos y presenta «se necesita más open human data» como conclusión de toda la clase; también recuerda que StableVicuna ya hizo funcionar PPO en 2023, aunque muy pocos equipos podían lograrlo.}

\subsection{Resumen de la sección}

Las variables decisivas de la primera ola de modelos instruct abiertos fueron, en orden: base weights disponibles, pipeline de instruction sintéticas, prompt de chat reales, datos humanos abiertos, licencias de publicación y un RLHF ejecutable. Los nombres de los modelos cambian sin cesar, pero la distribución de los datos y las interfaces reproducibles son una explicación más estable. La siguiente sección examinará cómo QLoRA redujo aún más la barrera de participación y por qué un fine-tuning de baja memoria no equivale automáticamente a un RLHF de bajo costo.
